%% kapa: exam=1 round=29 year=2018 period=1 form=A subject=civil_law subject_ko=민법 seq=28 number=28
%% source: Q-Net 감정평가사 기출문제 (cst00309) artlSeq=5202808 file=extracted/1008/5202808/2018년 제29회 감정평가사 1차시험 기출문제/2018년 제29회 감정평가사 1차시험 기출문제/1교시 A형.pdf
%% conversion: kordoc 4.24.51 -> tools/gen1.py (auto); review=proofread-sonnet; flags=none
\begin{question}{28}
소유권에 기한 물권적 청구권에 관한 설명으로 옳지 않은 것은? (다툼이 있으면 판례에 따름)
\begin{choices}
  \item 아직 건물의 소유권을 취득하지 못한 건물매수인은 그 건물의 불법점거자에 대하여 직접 건물의 명도청구를 할 수 없다.
  \item 소유물반환청구권의 상대방인 점유자가 그 물건을 점유할 권리가 있는 때에는 반환을 거부할 수 있다.
  \item 토지의 점유자가 점유취득시효를 완성한 경우에도 토지소유자는 그 토지의 인도를 청구할 수 있다.
  \item 소유권에 기한 물권적 청구권은 소멸시효의 대상이 되지 않는다.
  \item 소유물방해예방청구권에서 관념적인 방해의 가능성만으로는 방해의 염려가 있다고 할 수 없다.
\end{choices}
\end{question}
\answer{3}
